% =====================================================================
\chapter{Polinomio de Taylor}
\label{tema:6}
% =====================================================================

La recta tangente es la mejor aproximación de una función por un polinomio de grado~1 cerca de un punto. El polinomio de Taylor generaliza esa idea a cualquier grado: es el polinomio que tiene en el punto las mismas derivadas que la función hasta un orden dado. El teorema de Taylor con resto de Lagrange dice cuánto se equivoca esa aproximación, lo que permite calcular valores de funciones con un error garantizado, calcular límites y clasificar puntos críticos. Se siguen el capítulo~20 de Spivak~\cite{spivak}, el capítulo~7 de Apostol~\cite{apostol}, la §5.3 de los apuntes de \emph{Analysis~I}~\cite{cambA1}, la §6.3 del segundo volumen de OpenStax~\cite{osc2} y las sesiones~23--26 del curso 18.01SC del MIT~\cite{mit1801}.

\section{El polinomio de Taylor}

\begin{definicion}[Polinomio de Taylor]\label{def:taylor}
Sea $f$ derivable $n$ veces en $a$. El \emph{polinomio de Taylor} de orden $n$ de $f$ en $a$ es
\[
P_{n,a}f(x)=\sum_{k=0}^{n}\frac{f^{(k)}(a)}{k!}(x-a)^{k}=f(a)+f'(a)(x-a)+\frac{f''(a)}{2}(x-a)^{2}+\dots+\frac{f^{(n)}(a)}{n!}(x-a)^{n},
\]
y el \emph{resto} es $R_{n,a}f(x)=f(x)-P_{n,a}f(x)$. Si $a=0$ se habla de polinomio de \emph{Maclaurin}.
\end{definicion}

\begin{proposicion}[Caracterización]\label{prop:taylor-unico}
$P=P_{n,a}f$ es el único polinomio de grado a lo sumo $n$ que cumple $P^{(k)}(a)=f^{(k)}(a)$ para $k=0,1,\dots,n$.
\end{proposicion}

\addtocounter{definicion}{1}%

\section{El teorema de Taylor}

\begin{teorema}[Taylor con resto de Lagrange]\label{teo:taylor}
Sea $f$ derivable $n+1$ veces en un intervalo abierto $I$ que contiene a $a$. Para cada $x\in I$ existe $\xi$ entre $a$ y $x$ tal que
\[
f(x)=\sum_{k=0}^{n}\frac{f^{(k)}(a)}{k!}(x-a)^{k}+\frac{f^{(n+1)}(\xi)}{(n+1)!}(x-a)^{n+1}.
\]
\end{teorema}

Para $n=0$ el teorema es exactamente el teorema del valor medio: $f(x)=f(a)+f'(\xi)(x-a)$. El punto $\xi$ es desconocido, pero basta conocer una cota de $f^{(n+1)}$ para acotar el error.

\begin{corolario}[Cota del error y forma de Peano]\label{cor:cota}
Si además $\abs{f^{(n+1)}(t)}\le M$ para todo $t$ entre $a$ y $x$, entonces
\[
\abs{R_{n,a}f(x)}\le\frac{M}{(n+1)!}\abs{x-a}^{n+1}.
\]
En particular, si $f^{(n+1)}$ está acotada cerca de $a$, entonces $R_{n,a}f(x)/(x-a)^{n}\to0$ cuando $x\to a$, lo que se escribe $f(x)=P_{n,a}f(x)+o\bigl((x-a)^{n}\bigr)$.
\end{corolario}

La notación $o(g(x))$ («o pequeña») designa cualquier función $h$ con $h(x)/g(x)\to0$; se cumplen $o(x^{n})+o(x^{n})=o(x^{n})$, $x^{m}o(x^{n})=o(x^{n+m})$ y $o(x^{m})=o(x^{n})$ si $m\ge n$, cuando $x\to0$.

\begin{table}[h]
\centering\small
\caption{Desarrollos de Maclaurin más usados, cuando $x\to0$.}\label{tab:maclaurin}
\begin{tabular}{@{}ll@{}}
\toprule
$e^{x}=1+x+\dfrac{x^{2}}{2}+\dfrac{x^{3}}{6}+o(x^{3})$ & $\ln(1+x)=x-\dfrac{x^{2}}{2}+\dfrac{x^{3}}{3}+o(x^{3})$ \\[2mm]
$\sin x=x-\dfrac{x^{3}}{6}+o(x^{4})$ & $\cos x=1-\dfrac{x^{2}}{2}+\dfrac{x^{4}}{24}+o(x^{5})$ \\[2mm]
$(1+x)^{\alpha}=1+\alpha x+\dfrac{\alpha(\alpha-1)}{2}x^{2}+o(x^{2})$ & $\dfrac{1}{1-x}=1+x+x^{2}+x^{3}+o(x^{3})$ \\[2mm]
$\tanh x=x-\dfrac{x^{3}}{3}+o(x^{4})$ & $\sigma(x)=\dfrac12+\dfrac{x}{4}-\dfrac{x^{3}}{48}+o(x^{4})$ \\
\bottomrule
\end{tabular}
\end{table}

En la tabla, $\sin x=x-x^{3}/6+o(x^{4})$ (y no solo $o(x^{3})$) porque el término de grado~4 es nulo: $P_{3,0}\sin=P_{4,0}\sin$. Lo mismo ocurre con $\cos$, $\tanh$ y $\sigma-\tfrac12$, que son pares o impares.

\addtocounter{definicion}{1}%

\section{Aproximación con error controlado}

\addtocounter{definicion}{1}%

\addtocounter{definicion}{1}%

\section{Aplicaciones a límites y a extremos}

\addtocounter{definicion}{1}%

\begin{proposicion}[Criterio de la derivada $n$-ésima]\label{prop:criterio-n}
Sea $f$ de clase $C^{n}$ en un entorno de $a$, con $f'(a)=\dots=f^{(n-1)}(a)=0$ y $f^{(n)}(a)\ne0$, $n\ge2$. Si $n$ es par, $f$ tiene en $a$ un mínimo local estricto si $f^{(n)}(a)>0$ y un máximo local estricto si $f^{(n)}(a)<0$. Si $n$ es impar, $f$ no tiene extremo local en $a$.
\end{proposicion}

\addtocounter{definicion}{1}%

\addtocounter{definicion}{1}%

